
%%==================================================================
\section{Circuits and tensor networks}\label{sec:networks}
%%==================================================================

\subsection{The circuit classes}
\begin{definition}\label{def:gates}
A \emph{circuit} $C$ with $m$ inputs consists of input variables
$x_1,\dots,x_m$ and gates $g_1,\dots,g_G$. Each gate $g$ has an ordered list of
$f(g)\ge0$ predecessors, each an input or an earlier gate, and one gate is the
output. The \emph{wires} of $C$ are the pairs (predecessor, gate), counted with
multiplicity, so their number is $\wires(C)=\sum_gf(g)$, and $s(C)=m+\wires(C)$.
For integers $F,\beta\ge1$ the class $\cK(F,\beta)$ consists of the circuits in
which every gate $g$, with inputs $y=(y_1,\dots,y_f)$, is of one of the
following kinds.
\begin{enumerate}[label=\textup{(\alph*)},leftmargin=2.2em]
\item \emph{Table gate}: $f\le F$, and $g(y)=\tau(y)$ for a function
$\tau:\{0,1\}^f\to\{0,1\}$ given by its truth table of $2^f$ bits.
\item \emph{Weighted gate}: integers $w_1,\dots,w_f$ with $|w_i|\le2^\beta$ are
given, $v(y)=\sum_iw_iy_i$, and $g(y)=\varphi(v(y))$, where $\varphi$ is one of
\[
[v\ge\theta],\qquad [v=\theta],\qquad [v\equiv\theta\bmod\mu]\ \ (1\le\mu\le2^\beta),
\]
for a given integer $\theta$, or, when $w_1=\dots=w_f=1$, an arbitrary function
$\varphi:\{0,\dots,f\}\to\{0,1\}$ given by its table of $f+1$ bits.
\end{enumerate}
The \emph{description} of $C$ lists $m$, the gates in order with their kinds,
predecessor lists, tables, weights, thresholds and moduli in binary, and the
output gate; $L$ denotes its length.
\end{definition}

A threshold $\theta$ outside $[-f2^\beta,f2^\beta+1]$ can be replaced by the
nearest end of this interval without changing the gate, and $\theta$ in a
congruence can be reduced modulo $\mu$, so we may and do assume that all
numbers in a description have $O(\beta+\log f)$ bits. A table gate of fan-in
zero is a constant.

\begin{example}\label{ex:gates}
Every gate of fan-in at most two is a table gate. An AND or OR of any fan-in is
the exact gate $[\sum_iy_i=f]$ or the threshold $[\sum_iy_i\ge1]$; a negation
is $[{-y_1}\ge0]$; parity and $\mathrm{MOD}_\mu$ gates are congruences; majority
and every symmetric function are unit-weight gates with a table on
$\{0,\dots,f\}$. A linear threshold function whose minimal integer weights
exceed $2^\beta$ is not a weighted gate of $\cK(F,\beta)$, but it is a table gate
when $f\le F$.
\end{example}

\subsection{Tensor networks}
We use the definitions of \cite[Section D.1]{OAI26}. A \emph{Boolean tensor
network} over a commutative ring $R$ is a finite multigraph $\mathcal G$ in
which every edge carries one bit. Each vertex $v$ has an ordered list $\delta(v)$
of incident edges, its \emph{ports}, and a tensor $T_v:\{0,1\}^{\delta(v)}\to R$.
Parallel edges occupy distinct ports. The \emph{contraction} of the network is
\[
Z=\sum_{x\in\{0,1\}^{E(\mathcal G)}}\ \prod_{v\in V(\mathcal G)}T_v\bigl(x|_{\delta(v)}\bigr).
\]
The size of the network is its number of edges, and the network is
\emph{bipartite} if $\mathcal G$ is. For integers $P_0,\ell\ge1$ put
$B_{P_0,\ell}=(\Z/2^{P_0}\Z)[V]/(V^\ell)$.

\begin{theorem}[\oai{}, {\cite[Theorem D.1]{OAI26}}]\label{thm:contraction}
Fix positive constants $C,A,b$. Let $s\ge2$ be an integer and let $L$ be the
total length of the graph and tensor descriptions of a bipartite tensor network
with at most $Cs$ edges, maximum degree at most $\lceil s^{9/10}\rceil$, and a
number of vertices bounded by a fixed polynomial in $s$, whose entries lie in
$B_{P_0,\ell}$ with $1\le P_0,\ell\le Cs$. Suppose that every entry can be
computed from the description in time $(s+L)^b2^{As^{9/10}}$ and space
$(s+L)^b$. Then the contraction can be computed deterministically in time
$(s+L)^{O(1)}2^{O(s/\sqrt{\log s})}$ and space $(s+L)^{O(1)}$, with constants
depending only on $C,A,b$.
\end{theorem}

The proof is reproduced in Appendix~\ref{app:contraction}. We use only the case
$\ell=1$, in which $B_{P_0,1}=\Z/2^{P_0}\Z$.

\begin{figure}[t]
\centering
\includegraphics[width=0.66\textwidth]{fig_network.png}
\caption{The network of a small circuit (Section~\ref{sec:counting}). Gate
vertices (blue) carry the tables $[z=g(y)]$; each signal is distributed to its
ports by a tree of equality vertices (green), here a single vertex; a pin $(0,1)$ sits on the output;
every edge is then subdivided by a binary equality vertex (white), which makes
the network bipartite.}\label{fig:network}
\end{figure}

%%==================================================================
\section{Counting satisfying inputs}\label{sec:counting}
%%==================================================================

\begin{theorem}\label{thm:counting}
There are absolute constants $A_0,b_0$ and a deterministic algorithm that,
given the description of length $L$ of a circuit $C\in\cK(s^{4/5},s^{2/3})$ with
$s=s(C)\ge2$, computes $\#C=|\{x\in\{0,1\}^m:C(x)=1\}|$ in time
$(s+L)^{b_0}2^{A_0s/\sqrt{\log s}}$ and space $(s+L)^{b_0}$.
\end{theorem}

Here and below $s^{4/5}$ and $s^{2/3}$ stand for their integer parts. Since a
table gate has at most $s^{4/5}$ inputs, $L\le s^{O(1)}2^{s^{4/5}}$, and the time
is $2^{O(s/\sqrt{\log s})}$. The proof occupies the rest of this section.

\subsection{Normalisation}
A gate from which the output cannot be reached may be deleted; this removes its
wires and does not change the function. Constant gates of the same value may
be merged, which leaves the wires unchanged. We therefore assume that every
gate reaches the output and that there are at most two constant gates. Every
other gate has a wire, so
\begin{equation}\label{eq:gates}
G\le\wires(C)+2.
\end{equation}
Let $u$ be the number of inputs that are predecessors of no gate and that are
not the output. Then $\#C=2^u\,\#C'$, where $C'$ is $C$ with those inputs
removed, and $s(C')\le s(C)$. We assume $u=0$ from now on, restoring the factor
at the end.

\subsection{The network $\cN(C)$}
Put $K=\lfloor s^{4/5}\rfloor$, $\beta=\lfloor s^{2/3}\rfloor$ and
$b=\beta+\lceil\log K\rceil+2$. A gate is \emph{small} if it is a table gate or a
weighted gate of fan-in at most $K$, and \emph{big} otherwise. A \emph{signal}
is an input or a gate. For a signal $\sigma$ let $d_\sigma$ be its number of
\emph{occurrences}: one for its own value if $\sigma$ is a gate, one for each
wire leaving it, and one more if it is the output.

\smallskip\noindent\emph{Gate vertices.} A small gate $g$ with inputs
$y_1,\dots,y_f$ becomes a vertex $v_g$ with $f+1$ ports $(y_1,\dots,y_f,z)$ and
tensor $T(y,z)=[z=g(y)]$.

\smallskip\noindent\emph{Blocks.} A big weighted gate $g$ of fan-in $f>K$ has its
inputs split, in order, into $k=\lceil f/K\rceil$ blocks $I_1,\dots,I_k$ of at most
$K$ inputs each. For each block there is a vertex $v_{g,j}$ with ports
$(y_i)_{i\in I_j}$ and $b$ further ports $e_j=(e_{j,1},\dots,e_{j,b})$, and tensor
\[
T(y,e_j)=\bigl[\,\mathrm{val}(e_j)=P_j(y)+K2^\beta\,\bigr],\qquad
P_j(y)=\sum_{i\in I_j}w_iy_i,
\]
where $\mathrm{val}(e)=\sum_te_t2^{t-1}$. Since $|P_j|\le K2^\beta$, the integer
$P_j+K2^\beta$ lies in $[0,2K2^\beta]$ and has a unique binary expansion with
$b$ bits. A root vertex $r_g$ has the $kb$ ports $e_1,\dots,e_k$ and one port $z$,
with tensor
\[
T(e_1,\dots,e_k,z)=\Bigl[\,z=\varphi\Bigl(\sum_{j=1}^k\bigl(\mathrm{val}(e_j)-K2^\beta\bigr)\Bigr)\Bigr].
\]
The ports $e_j$ of $v_{g,j}$ and of $r_g$ are joined by $b$ edges, bit by bit.
For a big gate, the occurrence of a wire into $g$ is the corresponding input
port of a block vertex, and the occurrence of the value of $g$ is the port $z$
of the root.

\smallskip\noindent\emph{Equality trees and the pin.} Every occurrence of a
signal is a port of some vertex: an input port $y_i$ for a wire, the port $z$
for the value of a gate, and, for the output, the single port of a pin vertex
with tensor $(T(0),T(1))=(0,1)$. The $d_\sigma$ occurrences of $\sigma$ are
joined by an equality tree $Q_\sigma$: for $d_\sigma=1$ one edge to a vertex with
unary tensor $(1,1)$; for $d_\sigma=2$ two edges to a vertex with tensor $[a=a']$;
for $d_\sigma\ge3$ a tree with $d_\sigma$ leaf edges and $d_\sigma-2$ internal
vertices of degree three, each with tensor $[a=a'=a'']$. In every case each
common value on the leaf edges extends to exactly one assignment of the internal
edges with all tensors equal to one, and every other assignment has product
zero.

\smallskip\noindent\emph{Subdivision.} Finally every edge is subdivided once by a
vertex of degree two with tensor $[a=a']$. The vertices present before
subdivision form one side and the subdivision vertices the other, so $\cN(C)$
is bipartite. All entries are $0$ or $1$, read in $\Z/2^{m+1}\Z$.

Figures~\ref{fig:network} and~\ref{fig:blocks} show the construction.

\begin{figure}[t]
\centering
\includegraphics[width=0.92\textwidth]{fig_blocks.png}
\caption{A big weighted gate. Its $f$ inputs are cut into blocks of at most $K$
inputs; the block vertex $v_{g,j}$ writes the shifted partial sum
$P_j+K2^\beta$ on $b$ edges, and the root $r_g$ adds the decoded values and
applies $\varphi$. Every vertex then has degree $O(s^{13/15})$.}\label{fig:blocks}
\end{figure}

\begin{lemma}[Correctness]\label{lem:correct}
The contraction of $\cN(C)$ equals $\#C$ as an integer, and hence modulo
$2^{m+1}$.
\end{lemma}

\begin{proof}
Fix an assignment of the edges whose product of tensors is nonzero. The
subdivision vertices force both halves of each original edge to carry the same
bit. Each equality tree forces all occurrences of a signal $\sigma$ to carry a
common bit $a_\sigma$. At a small gate, $a_g=g(a_{\mathrm{pred}})$. At a big gate
the block tensors force $\mathrm{val}(e_j)=P_j(a)+K2^\beta$, so the sum read by
the root is $\sum_jP_j(a)=v(a)$ and $a_g=\varphi(v(a))$. The pin forces
$a_{\mathrm{out}}=1$. Conversely, an input $x$ with $C(x)=1$ determines the values
of all gates, hence all bits on all edges, and the product of the tensors at
that assignment is one. Distinct inputs give distinct assignments, and an input
with $C(x)=0$ gives none, because the pin would read zero. Since $u=0$, every
input occurs somewhere, so the nonzero terms are in bijection with the
satisfying inputs, and each is one.
\end{proof}

\begin{lemma}[Size]\label{lem:size}
For all $s$ larger than an absolute constant, $\cN(C)$ has at most $22s$ edges,
maximum degree at most $\lceil s^{9/10}\rceil$, and $O(s)$ vertices, and every
entry can be computed from the description of $C$ in time $(s+L)^{O(1)}$.
\end{lemma}

\begin{proof}
Let $W=\wires(C)$. The equality tree of $\sigma$ has at most $2d_\sigma$ edges, and
\[
\sum_\sigma d_\sigma=G+W+1,
\]
since each gate value, each wire and the output contribute one occurrence. A big
gate of fan-in $f$ has $k\le2f/K$ blocks and so $kb\le2fb/K$ edges between its
blocks and its root, and $b\le K$ for large $s$. Before subdivision the number
of edges is therefore at most $2(G+W+1)+2W\le6W+6\le11s$, using
\eqref{eq:gates} and $s\ge W$, $s\ge2$; subdivision doubles it. The number of
vertices is at most the number of edges plus the number of pins and unary
vertices, so it is $O(s)$.

Subdivision vertices have degree two and equality vertices degree at most
three. A small gate has degree $f+1\le K+1$, a block vertex at most $K+b$, and a
root
\[
kb+1\le\Bigl(\frac{s}{K}+1\Bigr)\bigl(s^{2/3}+\log s+3\bigr)+1=O(s^{1/5+2/3})=O(s^{13/15}).
\]
All of these are below $s^{9/10}$ for large $s$. An entry is a comparison of a
gate value with a table entry, a weighted sum of at most $s$ terms of $O(\beta+
\log s)$ bits, a binary expansion or a test of $\varphi$; each takes time
polynomial in $s+L$.
\end{proof}

\begin{proof}[Proof of Theorem~\ref{thm:counting}]
If $s$ is below the constant of Lemma~\ref{lem:size}, enumerate the $2^m$ inputs.
Otherwise normalise $C$, build $\cN(C)$ with its description of length
$(s+L)^{O(1)}$, and apply Theorem~\ref{thm:contraction} with $C=22$,
$P_0=m+1\le s$, $\ell=1$, any $A>0$ and the polynomial of Lemma~\ref{lem:size}.
By Lemma~\ref{lem:correct} the result is $\#C'\bmod2^{m+1}$. Multiply by $2^u$
modulo $2^{m+1}$; since $0\le\#C\le2^m$, the least nonnegative residue is $\#C$.
The bounds are those of Theorem~\ref{thm:contraction}.
\end{proof}

\begin{remark}\label{rem:space}
The algorithm uses space $(s+L)^{O(1)}$. A multitape Turing machine therefore
simulates it with a factor $(s+L)^{O(1)}$ in time, by scanning its workspace at
each access; this is the form used in Section~\ref{sec:lower}.
\end{remark}

\begin{corollary}\label{cor:classes}
For circuits whose gates are any mixture of gates of fan-in two, AND, OR and
parity gates of any fan-in, $\mathrm{MOD}_\mu$ gates with $\mu\le2^{s^{2/3}}$,
symmetric gates, and threshold gates with weights below $2^{s^{2/3}}$, the
number of satisfying inputs can be computed in time $L^{O(1)}2^{O(s/\sqrt{\log
s})}$, where $s$ is the number of inputs plus wires.
\end{corollary}

\begin{proof}
These circuits lie in $\cK(2,s^{2/3})$ by Example~\ref{ex:gates}, and here
$L=s^{O(1)}$ up to the bits of the weights, which are counted in $L$.
\end{proof}

%%==================================================================
\section{Zero-one programs}\label{sec:ilp}
%%==================================================================

\begin{corollary}[Theorem C]\label{cor:ilp}
Let $A\in\Z^{r\times n}$, $b\in\Z^r$ and $c\in\Z^n$, and put $\sigma=n+\nu(A)$,
where $\nu(A)$ is the number of nonzero entries of $A$. Suppose that every
entry of $A$ and $c$ has absolute value at most $2^{\sigma^{2/3}}$, and let $L$ be
the length of the input. Then, deterministically in time
$L^{O(1)}2^{O(\sigma/\sqrt{\log\sigma})}$ and polynomial space, one can
\begin{enumerate}[label=\textup{(\roman*)},leftmargin=2.2em]
\item compute the number of $x\in\{0,1\}^n$ with $Ax\le b$;
\item decide whether such $x$ exists and, if so, compute $\max\{c\cdot x:Ax\le
b,\ x\in\{0,1\}^n\}$ and a point attaining it.
\end{enumerate}
\end{corollary}

\begin{proof}
A row with no nonzero entry is the constant condition $0\le b_j$; if one fails,
the answer is zero, and otherwise drop it. So every row has a nonzero entry and
$r\le\nu(A)$. Row $j$ becomes the weighted gate
$h_j=[\sum_k(-a_{jk})x_k\ge-b_j]$, with $\nu_j$ wires, and their conjunction is the
exact gate $[\sum_jh_j=r]$ with $r$ wires. The circuit has $s=n+\nu(A)+r\le2\sigma$
and weights at most $2^{\sigma^{2/3}}\le2^{s^{2/3}}$, so it lies in
$\cK(2,s^{2/3})$ and Theorem~\ref{thm:counting} gives (i).

For (ii) add the gate $[c\cdot x\ge t]$, with $n$ wires, to the conjunction; the
circuit now has $s\le3\sigma$, and its weights are still at most
$2^{s^{2/3}}$. The value $\max c\cdot x$ lies in $[-n2^{\sigma^{2/3}},
n2^{\sigma^{2/3}}]$, so binary search over $t$ finds it with
$O(\sigma^{2/3}+\log n)$ counting calls. A maximiser is found by fixing
$x_1,x_2,\dots$ in turn, replacing $b$ by $b-a_kx_k$ and keeping a value for which
the count with $c\cdot x\ge t_{\max}$ stays positive; this uses $n$ further calls,
each on a smaller program. The calls are made one after another, so the space
stays polynomial.
\end{proof}

When $\nu(A)=O(n)$ the bound is $2^{O(n/\sqrt{\log n})}$. Problems with $O(n)$
constraints of bounded support, such as exact cover by sets of bounded size or
knapsack with conflict constraints along the edges of a sparse graph, are of
this form, with weights of up to $\sigma^{2/3}$ bits in the knapsack row.
